\documentclass[11pt]{article}

\usepackage[margin=1in]{geometry}
\usepackage{amsmath,amssymb}
\usepackage{enumitem}
\usepackage[hidelinks]{hyperref}
\usepackage{xcolor}

\setlist[itemize]{leftmargin=1.2em,itemsep=0.25em,topsep=0.25em}
\setlist[enumerate]{leftmargin=1.4em,itemsep=0.25em,topsep=0.25em}
\setlength{\parindent}{0pt}
\setlength{\parskip}{0.65em}

\title{Statistics-Dynamics Reframe for Puzzle RL Theory}
\author{Murphy}
\date{2026-04-24}

\begin{document}
\maketitle

\section*{Correction}

The scientific object should not be ``try several train/test-time scaling methods
on a puzzle-pretrained transformer and see what works.'' The corrected object is
to use a standard puzzle-training pipeline as a controlled dynamical system and
discover statistics whose trajectories distinguish imitation-like learning from
RL-like exploration or selection.

For branch $m$, checkpoint $t$, fixed prompt set $P$, and rollouts
$\tau \sim \pi_{m,t}(\cdot \mid x)$, the unit of analysis is a time series
\[
\Phi_m(t) =
\operatorname{observable}(\{\tau_{x,i,t}\}, \text{logits}, \text{gradients},
\text{activations}, \text{optimizer state}).
\]
The deliverable is a statistic or mathematical object with a reproducible,
interpretable shape, not a final-score ranking.

\section*{What the Coffee Automaton contributes}

The relevant lesson from Aaronson, Carroll, and Ouellette's Coffee Automaton
paper (arXiv:1405.6903) is methodological.

\begin{enumerate}
\item Use a simple controlled dynamics. The paper studies a small stochastic
cellular automaton rather than a realistic fluid simulator. The puzzle analogue
is a fixed generator, verifier, tokenizer, architecture, optimizer family, and
checkpoint schedule.
\item Separate entropy-like variables from complexity-like variables. Entropy
can rise monotonically while apparent complexity rises and then falls. In the
puzzle setting, loss, reward, and pass@1 may be necessary but not mechanistic.
\item Tie coarse-graining to causal structure. The automaton uses local spatial
coarse-graining. Maze and Sudoku statistics should be computed over parsed
puzzle states, legal actions, constraints, and rollout prefixes, not raw token
strings alone.
\item Validate against null dynamics. The paper's interacting versus
non-interacting comparison is the template. Puzzle nulls should preserve local
surface structure while breaking global solution structure.
\item Assume the first metric is artifact-prone. The paper finds and repairs a
threshold artifact. Puzzle analogues include compression serialization, cluster
thresholds, rollout budget, temperature, max length, action masks, and probe
shortcuts.
\end{enumerate}

\section*{Standard pipeline vessel}

The first pass should keep the training pipeline deliberately plain:

\begin{enumerate}
\item Pretrain on random valid puzzle trajectories.
\item Branch to SFT on solver traces.
\item Branch to REINFORCE-style RL with binary task reward.
\item Run AdamW and Muon as robustness branches for the statistics.
\end{enumerate}

Architecture, tokenizer/action format, task generator, verifier, probe prompts,
rollout count, temperature, max length, and checkpoint cadence should stay
fixed. Search-plus-distill, SePT-style self-training, on-policy supervision,
dense credit, KL/zero-KL RLVR, and exploration-first variants should be treated
as later perturbations once candidate observables exist.

\section*{Candidate objects}

\textbf{Scalar and optimizer dynamics.}
Track supervised loss, random-valid NLL, binary reward, legal-action entropy, KL
to pretrain, gradient norms, update norms, SFT/RL gradient alignment, AdamW/Muon
update effective rank, and Fisher or sharpness proxies. The interesting pattern
would be a nonmonotone expansion/contraction object, for example update-rank or
legal-action entropy, while reward or loss changes more monotonically.

\textbf{Behavior and process dynamics.}
Track pass@1, pass@k, trace-validity, final validity, optimality, invalid-action
rate by step, legal branching factor, error correction, and time-to-first-success.
A sharp distinction would be RL showing a first-success or invalid-action phase
transition while SFT shows smoother trace imitation.

\textbf{Exploration and support dynamics.}
Track state/action coverage, first-visit rates, occupancy entropy, trajectory
cluster birth/death, length-normalized base-pretrain logprob of successful
rollouts, and the fraction of successes already found by large-k pretrain
sampling. This is the most central family for support expansion versus
sharpening.

\textbf{Representation dynamics.}
Track hidden-state effective rank, probe margins for state/action/future success,
compression of quantized activations, policy-field apparent complexity, and
light-cone-like prefix predictability:
\[
\widehat I \approx \operatorname{CE}(z_{>t}) -
\operatorname{CE}(z_{>t} \mid h_{\leq t}),
\]
where $z_{>t}$ is a future event such as eventual success, next invalid action,
or return to a shortest-path frontier. These are conjectural and should come
after behavior logging is stable.

\textbf{Algorithmic complexity proxies.}
Track compressed description length of rollout sets, MDL of trajectory clusters,
minimal automata or grammars over strategy traces, and normalized compression
distance between checkpoints. These are closest in spirit to apparent complexity
but easiest to corrupt with symbolization choices.

\section*{First statistics to implement}

Start with four high-value statistics:

\begin{enumerate}
\item \textbf{Trace-validity phase curve:} legal-transition rate, correct
termination, invalid-action rate by step, shortest-path regret.
\item \textbf{Support occupancy entropy and first-visit rate:} coverage over
parsed puzzle states/actions under a fixed rollout budget.
\item \textbf{Trajectory-cluster birth/death plus base logprob:} clusters over
parsed trajectories, not raw strings, with length-normalized logprob under the
pretrain checkpoint.
\item \textbf{Policy-field apparent complexity:} evaluate next-action
distributions over a fixed grid of coarse puzzle states, quantize, puzzle-causal
coarse-grain, and compress.
\end{enumerate}

Then add one representation diagnostic, such as hidden-state effective rank or a
prefix probe for future success.

\section*{Nulls and artifact checks}

Use random legal-walk policies, target-shuffled prompts, reward permutation,
local-validity-only traces, corrupted solver traces, rollout-order shuffle,
compression algorithm sweeps, quantization sweeps, temperature sweeps, and
cluster-threshold sweeps. A statistic is not a candidate scientific object until
it survives these checks.

The desired first discovery is a statement of the form:

\begin{quote}
Policy-field apparent complexity rises and falls under RL but not SFT; the hump
coincides with cluster birth and first-success transitions; it disappears under
random-reward and local-validity nulls; and it is robust to compression and
quantization choices.
\end{quote}

\section*{Thesis}

The project is not trying to decide which post-training method wins on puzzles.
It is using a standard puzzle-training pipeline as a vessel for discovering
observables. Random-valid pretraining, solver-trace SFT, and binary-reward
REINFORCE give three controlled branches of the same dynamical system. The
scientific question is whether some statistic--support occupancy entropy,
trajectory-cluster birth/death, base-logprob of successes, policy-field apparent
complexity, representation rank, or prefix predictability--has a training-time
shape that distinguishes SFT-like imitation from RL-like exploration/selection.
A result matters only if the statistic is tied to puzzle causality, survives
null variants and artifact checks, and explains a mechanism that final reward
alone cannot show.

\end{document}
